% TOP_PROBLEM: {"problemId": "problem.cartan-hadamard-isoperimetric-conjecture", "canonicalTitle": "Cartan-Hadamard isoperimetric conjecture", "releaseRank": 274, "primaryDomain": "geometry_topology", "primaryDomainLabel": "Geometry & topology", "displayStatus": "open_with_solved_subcases", "releaseStatus": "open_with_solved_subcases"}
% !TeX program = lualatex
% ATTEMPT_STATUS: unresolved
\documentclass[11pt]{article}
\usepackage[margin=25mm]{geometry}
\usepackage{fontspec}
\setmainfont{Latin Modern Roman}
\usepackage{amsmath,amssymb,mathtools}
\usepackage{unicode-math}
\setmathfont{Latin Modern Math}
\usepackage{xurl,hyperref}
\hypersetup{colorlinks=true,urlcolor=blue}
\setcounter{secnumdepth}{0}
\setlength{\parindent}{0pt}
\setlength{\parskip}{5pt}
\setlength{\emergencystretch}{3em}
\begin{document}
\section*{\texorpdfstring{Cartan-Hadamard isoperimetric conjecture}{Cartan-Hadamard isoperimetric conjecture}}\label{274}
\subsection{Definitions and mathematical statement}
A Cartan--Hadamard manifold is a complete simply connected Riemannian manifold with sectional curvature at most zero. For a bounded measurable region $\Omega$ of finite perimeter and positive volume, write $V=\operatorname{vol}(\Omega)$ and $P(\Omega)$ for the total variation of its indicator function, agreeing with boundary area for smooth regions. If $\omega_n$ is the Euclidean unit-ball volume, the conjecture is
\[
 P(\Omega)\ge n\omega_n^{1/n}V^{(n-1)/n}.
\]
Equality is to occur only when the region, up to null sets, is isometric to a Euclidean ball in the source's rigidity sense. This does not force the entire ambient manifold to be Euclidean. The assertion is not limited to geodesic balls or smooth convex regions.
\par\smallskip\textit{Formulation and concept references:} \href{https://arxiv.org/html/2605.24638v1}{[S1]}.

\subsection{Short English statement}
Does nonpositive sectional curvature preserve the Euclidean lower bound for perimeter at fixed volume, with equality only for a Euclidean ball?
\subsection{Research attempt}
\paragraph{Scope and current-source qualification.}
This attempt was prepared by GPT-6 Astra Ultra on 27 September 2026.
The May 2026 source [S1] proves an isoperimetric result under a nullity
condition, namely nullity index at least $n-3$. Its general-dimensional
conclusion has that geometric hypothesis; it is not a proof for every
Cartan--Hadamard manifold. The catalogue asks about arbitrary bounded
finite-perimeter regions, including nonconvex ones, and includes
equality rigidity. Those requirements are retained.

The present work proves the sharp inequality and its rigidity for
geodesic balls in every Cartan--Hadamard manifold. It then tests why the
radial comparison used in that proof does not cover arbitrary regions.
The result is a restricted theorem, not the general conjecture.
Assume $n\ge2$ below. In dimension one, a complete simply connected
manifold is a line and a nonempty bounded finite-perimeter region has
at least two boundary points, with equality for an interval up to null sets.

\paragraph{Radial comparison with its sign checked.}
Fix $o\in M$. The Cartan--Hadamard theorem gives global geodesic polar
coordinates away from $o$, without cut or conjugate points. Write
\[
 d\operatorname{vol}=j(r,\theta)\,dr\,d\theta,
 \qquad A(r)=\int_{\mathbb S^{n-1}}j(r,\theta)\,d\theta,
 \qquad V(r)=\int_0^r A(s)\,ds.
\]
Here $d\theta$ is the standard unit-sphere measure, of total mass
$n\omega_n$. We claim
\[
 \partial_r\log j(r,\theta)\ge\frac{n-1}{r},
 \qquad j(r,\theta)\ge r^{n-1}.                         \tag{1}
\]
The geometric input can be seen directly from the index form.
For a transverse endpoint vector $v$ along a unit radial geodesic
$\gamma:[0,r]\to M$, let $J$ be its Jacobi field with $J(0)=0$
and $J(r)=v$. The second fundamental form of the distance sphere is
\[
 \operatorname{Hess}r(v,v)=
 \int_0^r\bigl(|D_sJ|^2-
       \langle R(J,\dot\gamma)\dot\gamma,J\rangle\bigr)\,ds.
\]
Sectional curvature at most zero makes the curvature contribution
nonnegative. Parallel-transporting to one tangent space and applying
Cauchy--Schwarz gives $\int_0^r|D_sJ|^2\,ds\ge|v|^2/r$.
Taking the transverse trace proves the first inequality in (1).
Since $j(r,\theta)/r^{n-1}\to1$ as $r\downarrow0$, integration gives
the second and also monotonicity of that ratio.
The sign is expansion relative to Euclidean space, not contraction.

\paragraph{A sharp inequality for every geodesic ball.}
For $0<s\le r$, monotonicity of $j(s,\theta)/s^{n-1}$ gives
\[
             A(s)\le A(r)(s/r)^{n-1},\qquad
             V(r)\le rA(r)/n.                           \tag{2}
\]
Also $A(r)\ge n\omega_n r^{n-1}$ by (1).
Combining the inequalities in the correct direction,
\[
 \begin{split}
 A(r)^n
 &=A(r)^{n-1}A(r)\\
 &\ge (nV(r)/r)^{n-1}\,n\omega_n r^{n-1}
   =n^n\omega_n V(r)^{n-1}.
 \end{split}
\]
Therefore
\[
                     A(r)\ge n\omega_n^{1/n}V(r)^{(n-1)/n}.
                                                               \tag{3}
\]
This proves the selected constant for all radii and centers within
the class of geodesic balls. Merely knowing both $A(r)$ and $V(r)$
are at least their Euclidean counterparts would not suffice;
the upper relation between $V(r)$ and $A(r)$ in (2) is essential.

\paragraph{Rigidity for the restricted theorem.}
Suppose equality holds in (3) for a positive radius. Equality must
hold in both factors used above, in particular
$A(r)=n\omega_n r^{n-1}$. Since
$j(r,\theta)\ge r^{n-1}$ continuously in $\theta$, equality of the
integrals forces equality for every direction. Monotonicity in (1)
then gives $j(s,\theta)=s^{n-1}$ for $0<s\le r$.

The shape operator $S_s$ of every radial sphere satisfies
$S_s\ge s^{-1}I$ by the index-form proof and has trace $(n-1)/s$.
Thus $S_s=s^{-1}I$. If $g_s$ denotes the angular metric in polar
coordinates, the relation $\partial_sg_s=2S_sg_s$ implies
$\partial_s(s^{-2}g_s)=0$. Its limit at zero is the round unit-sphere
metric, so
\[
                         g=ds^2+s^2g_{\mathbb S^{n-1}}
\]
throughout the ball. Hence the ball is isometric to a Euclidean ball.
This argument imposes no curvature condition outside it. It therefore
has exactly the local form of rigidity requested, within the class
we have actually proved.

\paragraph{A model-space check.}
In constant curvature $-\kappa^2$ with $\kappa>0$, put
$f(r)=\sinh(\kappa r)/\kappa$ and $c=n\omega_n$.
Then $A(r)=cf(r)^{n-1}$ and $V(r)=c\int_0^rf(s)^{n-1}\,ds$.
Since $f'(r)=\cosh(\kappa r)\ge1$,
\[
                  \int_0^rf(s)^{n-1}\,ds\le f(r)^n/n.
\]
Substitution gives (3) immediately, with strict inequality for $r>0$.
This independent check detects both the perimeter exponent and the
normalization of $\omega_n$.

There is also a local curvature check. The usual normal-coordinate
volume-density expansion gives, at a center $o$,
\[
 \begin{split}
 V(r)&=\omega_nr^n\left(1-
             \frac{\operatorname{Scal}(o)}{6(n+2)}r^2+O(r^3)\right),\\
 A(r)&=n\omega_nr^{n-1}\left(1-
             \frac{\operatorname{Scal}(o)}{6n}r^2+O(r^3)\right).
 \end{split}
\]
Consequently the ratio of the two sides of (3) is
\[
 1-\frac{\operatorname{Scal}(o)}{2n(n+2)}r^2+O(r^3).      \tag{4}
\]
Nonpositive sectional curvature implies nonpositive scalar curvature,
so the leading correction has the expected sign. This Taylor
calculation is a consistency check, not a substitute for the all-radius
proof or a proof about arbitrary shapes.

\paragraph{Why polar comparison does not settle arbitrary regions.}
Even for a smooth star-shaped region described by $0<r<R(\theta)$,
the volume is $\int\int_0^{R(\theta)}j(s,\theta)\,ds\,d\theta$,
whereas its boundary area involves angular derivatives of $R$ and
the full angular metric. Pointwise lower bounds for the radial density
increase both integrals and do not order their isoperimetric ratio.

The failed shortcut is already visible in Euclidean space. Ignoring
the angular-gradient term gives only
\[
 P\ge\int R^{n-1}\,d\theta,\qquad
 V=\frac1n\int R^n\,d\theta.
\]
On the sphere equipped with normalized probability measure, the
$L^{n-1}$ norm is at most the $L^n$ norm. Thus these two formulas alone
give the wrong direction for the desired perimeter lower bound.
Taking $R$ large on a small angular patch and small elsewhere makes
the ratio of those radial moments particularly small. Smoothing such
a profile produces a large angular-gradient contribution to the true
area; omitting it is the error, not a counterexample to isoperimetry.

A region may also fail to be star-shaped about every center, or have
several components and holes. Choosing a single center then does not
provide a one-valued radial boundary function at all. The radial-ball
argument has not addressed those cases.

\paragraph{The finite-perimeter extension requires a real premise.}
If the sharp inequality were established for every bounded smooth
region in a fixed manifold, it would extend to bounded finite-perimeter
regions by strict smooth approximation: choose smooth regions
$\Omega_j$ with characteristic functions converging in $L^1$,
volumes converging to $V$, and perimeters converging to $P(\Omega)$.
Passing to the limit in the smooth inequality then gives the same bound.
Local mollification in coordinate charts and coarea provide this
standard BV approximation in a relatively compact neighborhood.

Lower semicontinuity alone, $P(\Omega)\le\liminf_jP(\Omega_j)$,
would not transfer a lower bound from the approximants: it has the
wrong direction. Perimeter convergence is the crucial extra feature.
Moreover, approximating arbitrary regions by smooth ones does not
approximate them by geodesic balls, so it cannot extend our restricted
theorem (3) to the whole target. Equality rigidity for arbitrary BV
regions would require a further argument even after the inequality.

\paragraph{Verification and final status.}
The index form, the two inequality directions in (2), and the equality
case have been checked separately. The local script checks the scalar
coefficient in (4) algebraically and tests the hyperbolic formulas
numerically only as a normalization check. The infinite-dimensional
geometric proof is not certified by those samples.

\textbf{UNRESOLVED.} The sharp inequality and Euclidean-ball rigidity
are proved here for geodesic balls. The first missing step is a
perimeter-versus-volume estimate for arbitrary smooth regions without
the radial-ball assumption. Concrete next targets are an appropriate
rearrangement preserving the necessary area control, a global boundary
curvature argument beyond the nullity class of [S1], and equality
control under strict BV approximation. None has been supplied for
general Cartan--Hadamard manifolds.

\subsection{Sources}
\begin{itemize}
\item[S1]Mohammad Ghomi, Isoperimetric and total curvature inequalities in Cartan-Hadamard manifolds with nullity, arXiv:2605.24638v1, 2026.
\url{https://arxiv.org/html/2605.24638v1}.
\textit{Formulation source.}
\end{itemize}
\end{document}
